% Nastavení dokumentu
\documentclass[a4paper,12pt]{report}
\usepackage[utf8]{inputenc}

% Vlastní příkazy
\newcommand{\chp}[1]{\chapter*{#1}\addcontentsline{toc}{chapter}{#1}}
\newcommand{\schp}[1]{\section*{#1}\addcontentsline{toc}{section}{#1}}

% Nastavení Titulní strana: Nadpis, autor, datum
\title{Laboratorní cvičení čítače}
\author{Vojtěch Vašek}
\date{\today}

\begin{document}

% Tisk titulní strany, obsahu a vložení nové stránky
\maketitle
\tableofcontents
\newpage

\chp{Zadání a úvaha}
Zjistit přesnost čítače mikrokontroleru ATmega16 pomocí osciloskopu, porovnat vypočítanout hodnotu s reálnou a výsledky zobrazit v grafech.
Čítač je integrovaná periférie procesoru, která měří čas v podobě periody taktovacího signálu po průchodu el. obvodem dělení hodin. Tato periféri teké umí jednoduché počítání nahoru (přičte o jednu) a odčítání (odečte o jednu).
Čítač je jednodušeji obvod se schopností jednoduchých počtů složený z kombinační logiky a ragistrů.

\chp{Použité přístroje}
\schp{Seznam přístrojů}
\schp{Obrázky zapojení}
\chp{Postup práce}
\chp{Výpočty}
\chp{Tabulky}
\chp{Grafy}
\chp{Závěr}
\schp{Co jsem doufal, že se naučím a naučil}
\schp{Co jsem doufal, že se naučím a nenaučil}
\schp{Problémy s měřením}

\end{document}
